\section{Tổng quan về Đề tài}

\subsection{Giới thiệu Đề tài}

Trong bối cảnh chuyển đổi số diễn ra mạnh mẽ, các thư viện truyền thống đang dần chuyển mình thành các trung tâm học liệu trực tuyến hiện đại. Tuy nhiên, thách thức lớn nhất hiện nay không chỉ dừng lại ở việc số hóa danh mục sách mà còn ở khả năng quản lý và khai thác hiệu quả các tài nguyên đó. Các hệ thống quản lý thư viện hiện tại thường chỉ tập trung vào các tác vụ hành chính cơ bản, khiến việc tương tác giữa độc giả và nguồn tài liệu trở nên thụ động và rời rạc.

Đối với người quản lý thư viện, việc vận hành thủ công các danh mục tài nguyên khổng lồ—từ sách in, tạp chí đến các tài sản số—thường dẫn đến sự chậm trễ trong việc cập nhật trạng thái mượn trả và hồ sơ lưu trữ. Sự thiếu hụt các công cụ phân tích tự động tạo ra áp lực lớn lên công tác quản lý thành viên và theo dõi lưu thông, đặc biệt là khi dự báo nhu cầu đọc để tối ưu hóa việc bổ sung bộ sưu tập.

Đối với độc giả, việc tìm kiếm tài liệu trên các hệ thống cũ thường gặp khó khăn do cơ chế tìm kiếm dựa trên từ khóa cứng nhắc, thiếu tính linh hoạt và cá nhân hóa. Người dùng hiện đại mong muốn một không gian tương tác thông minh hơn, nơi họ có thể nhận được các gợi ý sách phù hợp với sở thích, lịch sử đọc và hành vi cá nhân mà không phải mất quá nhiều thời gian tìm kiếm thủ công. Đồng thời, nhu cầu được hỗ trợ tức thì qua các trợ lý ảo để giải đáp các thắc mắc về quy định hoặc vị trí giá sách đang trở nên thiết yếu.

Nhận thức được những yêu cầu cấp thiết đó, đề tài "Xây dựng Hệ thống Quản lý Thư viện Trực tuyến" được ra đời nhằm nghiên cứu, thiết kế và triển khai một nền tảng công nghệ toàn diện. Nền tảng này không chỉ giải quyết triệt để các bài toán quản lý tài nguyên và vận hành nghiệp vụ—như danh mục sách, tạp chí, tài sản số và quy trình lưu thông—mà còn tích hợp Trí tuệ nhân tạo (AI) để nâng cao hiệu quả khai thác bộ sưu tập. Trọng tâm của đề tài là tạo ra một hệ sinh thái tương tác thông minh dựa trên sự kết hợp giữa các mô hình quản lý hiện đại giúp tự động hóa vận hành và tập trung hóa quản lý tài nguyên, trải nghiệm người dùng cá nhân hóa thông qua các công cụ tìm kiếm thông minh và hệ thống gợi ý tự động dựa trên hành vi đọc thực tế, cùng sự hỗ trợ tự động từ Chatbot và trợ lý ảo để giải đáp thắc mắc và hướng dẫn vận hành cho người dùng mọi lúc, mọi nơi.

Bằng việc áp dụng các công nghệ phần mềm và thuật toán AI, hệ thống tập trung vào việc tối ưu hóa quy trình làm việc hành chính cho thủ thư, đồng thời cung cấp cho độc giả khả năng tìm kiếm nhanh chóng và tiếp cận tài liệu được cá nhân hóa theo nhu cầu.

\subsubsection{Động lực Chọn Đề tài}

Trong quá trình học tập và sử dụng thư viện tại môi trường đại học, nhóm chúng tôi nhận thấy nhiều hạn chế của các hệ thống thư viện hiện tại. Việc tìm kiếm tài liệu phụ thuộc quá nhiều vào từ khóa đơn giản, cho ra kết quả thường không sát với nhu cầu của người dùng. Hơn nữa, các hệ thống thiếu khả năng gợi ý tài liệu dựa trên sở thích hoặc lịch sử mượn, khiến độc giả mất nhiều thời gian để tìm kiếm tài nguyên phù hợp.

Bên cạnh đó, các hoạt động quản lý như theo dõi mượn/trả, cập nhật trạng thái mục mục hay hỗ trợ người dùng vẫn phụ thuộc nặng nề vào các thao tác thủ công. Điều này không chỉ làm giảm hiệu quả vận hành mà còn ảnh hưởng tiêu cực đến trải nghiệm của người dùng.

Nhận thức được thực tế này, nhóm chúng tôi đã chọn đề tài "Xây dựng Hệ thống Quản lý Thư viện Trực tuyến" với mục tiêu xây dựng một hệ thống hiện đại có khả năng tự động hóa và hỗ trợ thông minh, góp phần nâng cao hiệu quả quản lý và mức độ hài lòng của người dùng.

\subsection{Mục tiêu của Đề tài}

Mục tiêu chính của đồ án này là xây dựng một hệ thống quản lý thư viện trực tuyến dưới dạng Ứng dụng Web để hỗ trợ hiệu quả cho cả quản trị viên thư viện và độc giả trong việc khai thác, sử dụng tài nguyên học thuật.

Hệ thống được phát triển nhằm cung cấp một nền tảng quản lý tập trung cho các hoạt động thư viện, bao gồm quản lý danh mục, theo dõi vòng đời mượn trả và quản lý thành viên. Đồng thời, hệ thống hướng tới nâng cao trải nghiệm người dùng thông qua các công cụ tìm kiếm thông minh và cơ chế gợi ý dựa trên sở thích cũng như lịch sử đọc của người dùng.

Ngoài ra, đồ án hướng tới việc tích hợp các công nghệ trí tuệ nhân tạo để hỗ trợ người dùng trong quá trình tương tác với hệ thống. Các tính năng như chatbot hoặc trợ lý ảo được nghiên cứu và triển khai để đưa ra câu trả lời tức thì cho các câu hỏi thường gặp, hướng dẫn sử dụng hệ thống và hỗ trợ tra cứu tài liệu nhanh chóng.

Thông qua việc xây dựng hệ thống này, đồ án nỗ lực tạo ra một nền tảng thư viện số thông minh giúp tối ưu hóa quy trình làm việc hành chính cho các nhà quản lý thư viện, đồng thời mang lại trải nghiệm khám phá và truy cập tài nguyên thuận tiện, hiệu quả cho độc giả.

\subsubsection{Tổng quan về Hệ thống Đề xuất}

Hệ thống quản lý thư viện trực tuyến được xây dựng như một ứng dụng web, cho phép người dùng truy cập và tương tác thông qua các trình duyệt web tiêu chuẩn mà không cần cài đặt thêm phần mềm. Tổng thể, hệ thống bao gồm các thành phần cốt lõi: giao diện người dùng, động cơ xử lý logic nghiệp vụ phía máy chủ và cơ sở dữ liệu lưu trữ dữ liệu tả tài nguyên cùng hồ sơ người dùng.

Bên cạnh các thành phần cơ bản, hệ thống tích hợp các mô-đun trí tuệ nhân tạo để hỗ trợ các khả năng nâng cao như gợi ý tài liệu, tìm kiếm thông minh và chatbot hỗ trợ người dùng. Các mô-đun này hoạt động trên dữ liệu thu thập từ hành vi người dùng và dữ liệu tả tài nguyên của hệ thống để tạo ra các kết quả liên quan, mang tính cá nhân hóa.

\subsection{Lợi ích của Đề tài}

Đồ án mang lại những lợi ích thiết thực cho công tác quản lý và khai thác tài nguyên thư viện.

\subsubsection{Xây dựng Hệ thống Quản lý Thư viện Trực tuyến Hiện đại}

Đồ án tập trung phát triển hệ thống quản lý thư viện dựa trên web cho phép quản trị và truy cập tài nguyên thông qua môi trường trực tuyến. Hệ thống tích hợp các công nghệ xử lý dữ liệu và trí tuệ nhân tạo để nâng cao khả năng kiểm soát hành chính và hỗ trợ người dùng trong quá trình tìm kiếm, khai thác tài liệu.

\subsubsection{Hỗ trợ Hiệu quả cho Công tác Quản lý Thư viện}

Hệ thống cho phép quản lý các danh mục tài nguyên đa dạng, bao gồm sách in, tạp chí và các tài sản số. Các tính năng theo dõi trạng thái mượn-trả-bảo quản, tài khoản thành viên và phân tích sử dụng giúp nhân viên thư viện vận hành hiệu quả. Các thông tin thống kê từ hệ thống cũng hỗ trợ đưa ra quyết định liên quan đến việc bổ sung và điều chỉnh bộ sưu tập.

\subsubsection{Tối ưu hóa Việc Khai thác Tài nguyên của Độc giả}

Các khả năng tìm kiếm thông minh, động cơ gợi ý tự động và tính năng trợ lý ảo cho phép người dùng nhanh chóng tìm thấy các tài liệu phù hợp với nhu cầu học tập và nghiên cứu. Hệ thống cũng cá nhân hóa trải nghiệm người dùng dựa trên lịch sử đọc và thói quen tìm kiếm, tăng cường hiệu quả khai thác tài nguyên.

\subsubsection{Tăng cường Tương tác Giữa Người dùng và Hệ thống Thư viện}

Việc tích hợp chatbot hoặc trợ lý ảo đảm bảo người dùng nhận được sự hỗ trợ kịp thời khi tương tác với hệ thống. Các thắc mắc về chính sách mượn sách, vị trí tài liệu hoặc cách vận hành hệ thống được giải đáp ngay lập tức, nâng cao trải nghiệm tổng thể người dùng trong quá trình tương tác với thư viện.

\subsection{Phạm vi và Hạn chế của Đề tài}

Phạm vi của đồ án tập trung vào việc nghiên cứu và xây dựng hệ thống quản lý thư viện trực tuyến dưới dạng Ứng dụng Web. Các tác nhân trực tiếp của hệ thống bao gồm: Quản trị viên (Admin), Thủ thư (Librarian) và Độc giả (Reader).

\subsubsection{Các Nhóm Người dùng}

Độc giả (Reader): Người dùng chính của hệ thống, thực hiện tìm kiếm, tra cứu và mượn tài liệu. Độc giả cũng nhận được các gợi ý sách cá nhân hóa, quản lý lịch sử mượn và tương tác với hệ thống qua các tính năng tìm kiếm hoặc chatbot.

Thủ thư (Librarian): Chịu trách nhiệm quản lý các mục danh mục, xử lý các giao dịch mượn và trả, đồng thời cập nhật trạng thái tài nguyên. Thủ thư cũng sử dụng các công cụ quản trị để nhập dữ liệu và kiểm soát thông tin hệ thống.

Quản trị viên (Administrator/Admin): Quản lý toàn bộ hệ thống, bao gồm cấu hình tham số, quản lý tài khoản người dùng và giám sát tổng thể hệ thống.

Hệ thống quản lý các tài nguyên thư viện phổ biến như sách in, tạp chí in và các tài sản số ở dạng tệp điện tử như E-book hoặc tài liệu PDF. Các chức năng cốt lõi bao gồm quản lý danh mục, quản lý tài khoản thành viên, theo dõi trạng thái mượn-trả, hỗ trợ tìm kiếm tài liệu và tạo gợi ý cá nhân hóa.

Trong phạm vi đồ án, hệ thống được thiết kế cho các thư viện quy mô vừa và nhỏ, như thư viện khoa, thư viện trường học hoặc các trung tâm học liệu. Người dùng đích bao gồm sinh viên, giảng viên và các nhà nghiên cứu có nhu cầu truy cập tài liệu học thuật.

Các tính năng AI được triển khai ở mức độ nền tảng, được chi tiết hóa trong phần phạm vi ứng dụng trí tuệ nhân tạo dưới đây.

\subsubsection{Phạm vi Ứng dụng Trí tuệ Nhân tạo}

Trong phạm vi đồ án, các tính năng trí tuệ nhân tạo được triển khai ở mức độ nền tảng để minh họa tiềm năng tích hợp trong một hệ thống quản lý thư viện. Cụ thể, hệ thống sử dụng các mô hình và dịch vụ hiện có để triển khai các chức năng như thuật toán gợi ý theo hành vi, tìm kiếm bằng ngôn ngữ tự nhiên và chatbot Q\&A cho các thắc mắc thường gặp.

Do hạn chế về thời gian và tài nguyên, đồ án không tập trung vào việc xây dựng hoặc huấn luyện các mô hình AI từ đầu, mà chủ yếu tận dụng các công cụ và API hiện có để triển khai và tích hợp chúng vào hệ thống. Điều này đảm bảo tính khả thi của đồ án trong khi vẫn thể hiện được tiếp cận công nghệ hiện đại.

\subsubsection{Hạn chế của Đề tài}

Mặc dù hệ thống được thiết kế với các tính năng toàn diện, một số hạn chế nhất định vẫn tồn tại. Hệ thống hiện tại chỉ được phát triển dưới dạng ứng dụng web và chưa hỗ trợ các nền tảng di động gốc (native mobile). Các khả năng AI được triển khai ở mức độ cơ bản và chưa đạt đến sự tối ưu của các nền tảng thương mại lớn.

Ngoài ra, hệ thống chưa xử lý các kịch bản quy mô cực lớn liên quan đến hàng triệu bản ghi hoặc tích hợp sâu với các mạng lưới thư viện bên ngoài của bên thứ ba. Các khả năng nâng cao sẽ được xem xét cho các nghiên cứu tiếp theo và các phiên bản phần mềm sau.

\subsection{Ý nghĩa Tổng quan của Đề tài}

\subsubsection{Ý nghĩa Thực tiễn}

Hệ thống quản lý thư viện trực tuyến được phát triển trong đồ án này nâng cao hiệu quả vận hành cho các thư viện trong môi trường số. Việc tập trung hóa quản lý tài nguyên và người dùng giúp giảm bớt các thao tác thủ công, đồng thời đảm bảo tính chính xác của dữ liệu và sự đồng bộ theo thời gian thực.

Hơn nữa, hệ thống cải thiện đáng kể trải nghiệm đọc thông qua các công cụ tìm kiếm và gợi ý thông minh. Người dùng có thể nhanh chóng phát hiện các tài liệu phù hợp với nhu cầu học thuật mà không cần thực hiện các quy trình tìm kiếm phức tạp.

Ngoài ra, việc tích hợp chatbot hoặc trợ lý ảo giúp hỗ trợ tức thì cho người dùng, giảm tải công việc cho nhân viên thư viện trong việc giải đáp các thắc mắc thông thường.

\subsubsection{Ý nghĩa Khoa học}

Dưới góc độ nghiên cứu, đồ án làm rõ các phương pháp luận trong việc xây dựng và triển khai một hệ thống quản lý thư viện sử dụng các công nghệ phần mềm hiện đại. Thiết kế hệ thống bao gồm kỹ thuật kiến trúc phần mềm, thiết kế cơ sở dữ liệu và triển khai quản lý tài nguyên/người dùng.

Đồ án cũng khám phá việc tích hợp các mô hình trí tuệ nhân tạo vào hệ thống thông tin để cá nhân hóa trải nghiệm người dùng. Việc áp dụng các thuật toán gợi ý và chatbot hỗ trợ người dùng thể hiện tiềm năng của AI trong việc nâng cao hiệu quả khai thác tri thức.

Bằng cách kết hợp các công nghệ phát triển web, quản lý dữ liệu và trí tuệ nhân tạo, đồ án đề xuất một mô hình hệ thống thư viện thông minh có thể làm cơ sở tham khảo cho các nghiên cứu và triển khai thực tiễn trong tương lai.

Trong các chương tiếp theo, báo cáo tập trung vào phân tích yêu cầu hệ thống, thiết kế kiến trúc tổng thể và triển khai các tính năng cốt lõi để xây dựng một hệ thống quản lý thư viện hoàn chỉnh.